\section{Lecture 6: Aquatic pollution}

\subsection{Definitions}

\textbf{Pollution}
\begin{itemize}
	\item Eutrophication
	\item Organic matter
	\item Rubbish
	\item Microbial
	\item \textbf{Hazardous substances / contaminants / toxicants}
\end{itemize}

\subsection{Naturally occuring hazardous substances}

\begin{itemize}
	\item Heavy metals, PAHs, petroleum, hormones, ...
	\item essential and non-essential substances
\end{itemize}

Toxins -- substances produced by living organisms that are toxic to humans.

\subsection{Man-made hazardous substances}

\begin{itemize}
	\item 350000 registered
	\item 100000 on the market in Europe
\end{itemize}

\subsection{Planetary boundaries}

\textbf{2009}
\begin{itemize}
	\item 7 boundaries assessed
	\item 3 crossed
\end{itemize}

\textbf{2015}
\begin{itemize}
	\item 7 boundaries assessed
	\item 4 crossed
\end{itemize}

\textbf{2023}
\begin{itemize}
	\item 9 boundaries assessed
	\item 6 crossed
\end{itemize}

\subsection{Itai-itai (ouch-ouch) disease 1912}
\begin{itemize}
	\item cause: Cadmium (Cd) -- heavy metal
	\item in Japan
	\item from a mine
	\item the cadmium went into the rivers, and the rivers were used to
		water the rice plants
\end{itemize}

\subsection{Minamata disease 1956}
\begin{itemize}
	\item also Japan
	\item Mercury (Hg), also heavy metal
	\item Methyl mercury, a form of mercury transformed by bacteria
	\item again, originating from a factory
\end{itemize}

\subsection{Effects on wild-life 1960s}
\begin{itemize}
	\item A decline in eagles in Sweden and north of Baltic Sea
	\item Due to DDT toxicants
	\item Measured in DDT toxic equivalents
	\item Clear corelation between DDT toxic equivalents and egg shell
		thickness
\end{itemize}

\subsection{Biomagnification}
\begin{itemize}
	\item Concentration of DDT increases up the food chain
\end{itemize}

\subsection{Tri-butyl-tin (TBT)}

\begin{itemize}
	\item Oysters exposed to TBT have different shapes and become sterile
	\item The shape is very degenerate of course
\end{itemize}

\subsection{Oil spills}

\subsection{Sources, fate and status}
\begin{itemize}
	\item Input from diffuse and point sources
	\item Air/water exchange of Hz\footnote{hazardous substance}
	\item Wet + dry deposition of Hz
	\item Direct + indirect photolysis UV
	\item Chemical and biological transformation
	\item Sedimentation
	\item Chemical and biological transforming in the sedimentation
\end{itemize}

\subsection{Hazardous substances are everywhere}

\subsection{Sewage treatment plants}
Pain killers
\begin{itemize}
	\item Diklofenak
	\item Naproxen
	\item Ketoprofen
	\item Ibuprofen
\end{itemize}

PFAS
\begin{itemize}
	\item Recently contaminated drinking water in Skåne
\end{itemize}

\subsection{Assessing status -- Environmental quality standards}

\begin{itemize}
	\item Identify receptors and compartments at risk
	\item Collate and quality assess data
	\item Extrapolation
	\item Propose EQS
	\item Implement EQS
\end{itemize}

\subsection{Chemical risk -- fresh water}
Acute exposure is not so bad in Europe but chronic exposure is pretty bad.
Some regions have over 75\% chemical risk.

\subsection{Environmental impact -- ellpout larvea deformities}

\subsection{Environmental impact -- mussel hemolymph}

\subsection{Bottom-up -- Single species tests}
\begin{itemize}
	\item Standard test species and protocols
	\item Low biodiversity
	\item High repeatability and reproducibility
\end{itemize}

\subsection{Biodiversity -- why it matters}
\begin{itemize}
	\item Impact of sewage in Greenland
\end{itemize}

\subsection{Assessing reproductive success}

\subsection{Biodiversity and community ecotoxicology}

\textbf{Community}: an assemblage or associations of populations of two or more
different species, belonging to the same group, that occupy the same
location

\subsection{Top down -- a more ecological approach}

\begin{itemize}
	\item Secondary producers (SP)
	\item Primary producers (PP)
\end{itemize}

\subsection{Scrubbers -- closed loop}

\section{Lecture 6: plastic contamination}

\begin{itemize}
	\item 350000 chemicals registered for production globally
	\item over 16000 chemicals used in plastics
	\item 10726 chemicals without hazard data
	\item 4219 chemicals classified as hazardous
	\item 1192 chemicals less hazardous
\end{itemize}

\subsection{Impacts on biodiversity due to debris abundance}

\begin{itemize}
	\item Ingestion / entanglement
	\begin{itemize}
		\item Entanglement: starvation, drowning, limits foraging.
			Mostly due to "ghost fishing", but not only
		\item Ingestion: may lead to gut clogging, starvation.
			Accidental or confusion with prey
	\end{itemize}
	\item Habitat destruction
	\begin{itemize}
		\item Landfill, dumping
		\item Flooding
		\item Smothering:
		\begin{itemize}
			\item mangrove forests: weight and shading
			effect affects sensitive vegetation
		\end{itemize}
	\end{itemize}
\end{itemize}

\subsection{Impacts on biodiversity due to plastic-associated chemicals}

\subsection{Impacts of microplastics on biodiversity}

